



\section{Toric Encoded Circuits}

    \begin{definition}[Troic Encoded Circuit]
      Let $C$ be a quantum circuit. We say that $C$ is \textit{Toric encoded} if the following holds:
      \begin{enumerate}
          \item There is a Toric code $\mathcal{C}$ such that the qubits of $\mathcal{C}$ can be partitioned into registers $R_{1}, \ldots, R_{k}$, and each of the registers is encoded by $\mathcal{C}$.   
          \item The layers in $C$ with an even time coordinates are sweep-cycle.
        \end{enumerate} 
        For a set of faults $\mathcal{E}$ the gates in $C[\mathcal{E}]$ can be decomposed to cycles, of the following form:
        \begin{equation*}
          \begin{split}
            C =     D_{l}P_{l}\Lambda_{l},..,D_{1} P_{1} \Lambda_{1}
          \end{split}
        \end{equation*}
    \end{definition}

  \begin{definition}[Base graph $G_{o}$]
    Let $C$ be a quantum circuit, at depth $d$, with registers $R_{1}, \ldots, R_{k}$, all of which are encoded by Toric codes with the same parameters (i.e., dimension, length, etc.). Let's denote by $\mathcal{G}_{1}, .., \mathcal{G}_{k}$ the Toric complexes that induce the codes. We define the base graph of $C$, denoted by $G_{o} = (V_{o},E_{o})$ as follows: 
    \begin{enumerate}
      \item The vertices $V_{o}$ are tuples, where the first coordinates correspond to a vertex in the union over the vertices in $\mathcal{G}_{1}, \ldots, \mathcal{G}_{k}$, and the second coordinate is associated with a time. Namely:
        \begin{equation*}
          \begin{split}
            V_{o} \colonequals \bigsqcup_{i} V(\mathcal{G}_{i}) \times [d]
          \end{split}
        \end{equation*}
      \item The edges $E_{o}$ are derived from the original edges in $\mathcal{G}_{1}, \ldots, \mathcal{G}_{k}$, with the addition that any two vertices in preceding times $t$ and $t+1$ that support qubits (faces) such that $C$ acts non-trivially on those qubits at time $t$ are also connected by an edge. Namely:
        \begin{equation*}
          \begin{split}
            E_{o} & \colonequals \big\{ \{(v,t), (u,t^{\prime}) \} :  \text{ if }     \\ 
              & \text{either there exists } i \text{ such }  \{v,u\} \in \mathcal{G}_{i} \text{ and } t^{\prime} = t,t+1 \\
             & \text{or there exists } u_i = (g_{i},l_{i}) \in C \text{ such }  v,u \in l_{i,2} \text{ and } t = l_{i,1}, t^{\prime} = l_{i,1}+1  \big\} 
          \end{split}
        \end{equation*}
    \end{enumerate} 
\ctt{Here we should add that a vertex $v$ belongs to location $l$ if there is a qubit $f$ that $v$ belongs to its associated face. }

  \end{definition}

  \begin{definition}[Time Graph of $C(\mathcal{E})$]
    For a quantum circuit $C$ and set of faults $\mathcal{E}$, let us denote the base graph of $C[\mathcal{E}]$ by $G_{o}$. For each time $t$, denote the deviation induced by $\mathcal{E}$ at time $t$ by $\mathcal{D}_{t}$, so $\partial\mathcal{D}_{t}$ is the syndrome that the qubits in the deviation admit. We define the history graph $\mathcal{H} = ( V_{\mathcal{H}}, E_{\mathcal{H}})$ as follows: 
    \begin{enumerate}
      \item The vertices $V_{\mathcal{H}}$ are taken to be the vertices of $G_{o}$, namely tuples $(v,t) \in V_{o}$, that satisfy that $v$ belongs to a face associated with a check, which, at time $t$, belongs to the $\partial \mathcal{D}_{t}$. Namely:
        \begin{equation*}
          \begin{split}
            \big\{ v_{t} : v_{t}=(v,t) \in V_{o} \text{ and there exists a check } f \in C \text{ such } v \in f \in \partial \mathcal{D}_{t} \big\}
          \end{split}
        \end{equation*}
      \item The edges $E_{\mathcal{H}}$ is partitioned into two types $A_{\mathcal{H}}$ and $F_{\mathcal{H}}$, where $A_{\mathcal{H}}$ are the edges in $E_{o}$ restricted to $V_{\mathcal{H}}$ and $F_{\mathcal{H}}$ are defined to be the . Namely:
        \begin{equation*}
          \begin{split}
            A_{\mathcal{H}} &= \big\{ e = \{u,v\} \in E_{o} :  u,v \in V_{\mathcal{H}} \text{ and } \Time{v} = \Time{u}\pm 1 \big\} \\
            F_{\mathcal{H}} &= \big\{ e = \{u,v\} \in E_{o} : \text{ exists } w \in V_{\mathcal{H}} \text{ such } \{v,w\},\{u,w\} \in A_{\mathcal{H}} \big\} 
          \end{split}
        \end{equation*}
    \end{enumerate}
    By definition of the construction, $G_{\mathcal{H}}$ is a time-graph.
\end{definition}



\begin{definition}[Time Graph of $C(\mathcal{E})$]
  For a quantum circuit $C$ and set of faults $\mathcal{E}$, let us denote the base graph of $C[\mathcal{E}]$ by $G_{o}$. For each time $t$, denote the deviation induced by $\mathcal{E}$ at time $t$ by $\mathcal{D}_{t}$, so $\partial\mathcal{D}_{t}$ is the syndrome that the qubits in the deviation admit, denote by $\partial\mathcal{D}_{t}^{X}$ and by $\partial\mathcal{D}_{t}^{Z}$ the $X-Z$ decomposition of $\partial\mathcal{D}_{t}$. We define the history graph $\mathcal{H} = ( V_{\mathcal{H},X} \cup V_{\mathcal{H},Z}, E_{\mathcal{H}})$, in recursive manner as follow: 
  \begin{enumerate}
    \item For the empty circuit. The vertices in $V_{\mathcal{H}, X} $ are the vertices of $V_{o}|_{1}$ that support $\partial D_{t}^{X}$. Similarly, the vertices of $V_{\mathcal{H,Z}}$ are the vertices of $V_{o}|_{1}$ that support $\partial D_{t}^{Z}$. 
    \item Let $\mathcal{H}$ be the history graph of $C[\mathcal{E}]$ restricted to the first $d-1$ layers. Let $U$ be the $d$-layer, then define the history graph corresponding to the first $d$ layers of $C[\mathcal{E}]$ to be obtained by adding vertices and edges to $\mathcal{H}$ as follows:
      \begin{enumerate}
        \item If $U$ is a noise gate, namely $U \subset \mathcal{E}$, then denote by $\mathcal{E}_{t}^{X}$ and $\mathcal{E}_{t}^{Z}$ the $X,Z$ decomposition of $\mathcal{E}_{t}$. Then, for any vertex in $\mathcal{G}$
        \item Else, denote by $U_{1},..,U_{l}$ the gates in $U$. If $U_{1}$ is a $\mathbf{CX}^{\otimes n}$ gate acting on registers $R_{1}$ into $R_{2}$, then for each vertex $v \in V_{\mathcal{H},X}|_{R_{1}, t=d-1}$ add the counter parts of $v$ in $R_{1}$,$R_{2}$ at time $t=d+1$ into $V_{\mathcal{H} ,X}$ and connect them by arrows to $v$. Similarly, for any vertex $v \in V_{\mathcal{H},Z}|_{R_{2}, t=d-1}$ add the counter parts of $v$ in $R_{2}$,$R_{2}$ at time $t=d+1$ into $V_{\mathcal{H} ,X}$ and connect them by arrows to $v$.
      \end{enumerate}
  \end{enumerate}



\end{definition}







